\subsection{Clases}

\subsubsection{Regla: PascalCase y sustantivo singular}

\textbf{Regla:} El nombre de una clase se escribe en \texttt{PascalCase} y es un sustantivo o frase nominal en singular (Cwalina \& Abrams, 2008). Las clases de prueba se exceptúan porque siguen la convención específica del prefijo \texttt{Test} definida en este estándar.

\textbf{Código correcto:}
\begin{lstlisting}[style=csharpstyle]
public sealed class BoardTile
{
}
\end{lstlisting}

\textbf{Código incorrecto:}
\begin{lstlisting}[style=csharpstyle]
public sealed class BoardTiles
{
}
\end{lstlisting}

\textbf{Excepción:} Una clase \texttt{static} que solo agrupa valores predefinidos del mismo tipo se nombra en plural, como \texttt{Colors} o \texttt{Brushes} en .NET. Una clase estática no representa un objeto, porque no se pueden crear instancias de ella; representa el conjunto, y cada miembro se lee como un elemento de ese conjunto: \texttt{TileColors.Highlight}.

\textbf{Código correcto (clase estática):}
\begin{lstlisting}[style=csharpstyle]
public static class TileColors
{
    public static Color Highlight { get; } = new Color(0xF0, 0xC4, 0x88);

    public static Color Blocked { get; } = new Color(0xB3, 0x26, 0x1E);
}
\end{lstlisting}

\textbf{Código incorrecto (clase estática):}
\begin{lstlisting}[style=csharpstyle]
public static class TileColorList
{
    public static Color Highlight { get; } = new Color(0xF0, 0xC4, 0x88);

    public static Color Blocked { get; } = new Color(0xB3, 0x26, 0x1E);
}
\end{lstlisting}

\subsubsection{Regla: sin prefijos ni abreviaturas poco claras}

\textbf{Regla:} Sin prefijos que indiquen que un tipo es una clase (\texttt{CTile}, \texttt{ClsPlayer}) ni abreviaturas poco reconocidas (Cwalina \& Abrams, 2008).

\textbf{Código correcto:}
\begin{lstlisting}[style=csharpstyle]
public sealed class TurnManager
{
}
\end{lstlisting}

\textbf{Código incorrecto:}
\begin{lstlisting}[style=csharpstyle]
public sealed class ClsTrnMgr
{
}
\end{lstlisting}
